\chapter{Locks from the hardware up}
\label{ch:locks}

A lock turns a critical section into something atomic. But how do you
build a lock if every load and store can be interleaved with those of
other threads? This chapter answers it from the bottom: the requirements,
the one-CPU trick (disable interrupts), a classic software algorithm that
fails on real hardware, the atomic instructions that make locks possible,
and the choice between spinning and sleeping --- which is exactly the
difference between SWEB's \code{SpinLock} and \code{Mutex}.

\section{What a lock must guarantee}

\begin{description}
  \item[Mutual exclusion] at most one thread in the critical section.
  \item[Progress] if the lock is free and threads want it, one of them
    gets it (no deadlock among the contenders).
  \item[Bounded waiting] a waiting thread gets its turn after a bounded
    number of others (no starvation) --- ``fairness''.
\end{description}
and in practice: cheap when uncontended, and not burning CPU time while
waiting for long.

\section{Disabling interrupts}

On a single CPU, the only way another thread can run in the middle of
your critical section is a timer interrupt that makes the scheduler
switch. Turn interrupts off (\code{cli}), do the work, turn them back on
(\code{sti}) --- nothing can interrupt you. This is how kernels protected
data before multi-core, and SWEB still uses it for the scheduler's own
data (\code{Scheduler::schedule} runs with interrupts off). But:
\begin{itemize}
  \item it is privileged --- user code cannot do it;
  \item it does not stop the \emph{other} cores on a multi-core machine;
  \item with interrupts off, timer ticks and device interrupts are
    delayed; the section must be very short, and must not sleep.
\end{itemize}

\section{Peterson's algorithm --- and why it breaks}

With only loads and stores, two threads can still achieve mutual
exclusion on paper:

\begin{lstlisting}
int flag[2], turn;                 /* shared */
void lock(int me) {
  int other = 1 - me;
  flag[me] = 1;                    /* I want in           */
  turn = other;                    /* but you go first    */
  while (flag[other] && turn == other) {}   /* wait while you want in and it's your turn */
}
void unlock(int me) { flag[me] = 0; }
\end{lstlisting}

It is correct under \emph{sequential consistency} --- the assumption that
all memory operations happen in one global order that respects each
thread's program order. Real CPUs do not guarantee that. On x86, a store
goes into the core's \emph{store buffer} and becomes visible to the other
core later; a subsequent load from a \emph{different} address can
overtake it. Both threads store their \code{flag}, both load the other's
still-old flag = 0 from memory, and both enter (\module{04}'s
\code{peterson.c} shows the violation within seconds). A \emph{memory
fence} (\code{mfence}, \code{atomic\_thread\_fence(memory\_order\_seq\_cst)})
between the stores and the loads repairs it. The lesson: correct
synchronisation needs help from the hardware.

\section{Atomic instructions}

\begin{description}
  \item[Test-and-set / exchange] \code{old = xchg(\&lock, 1)} writes 1 and
    returns the previous value, atomically. If \code{old} was 0, you have
    the lock. (\code{xchg} with memory is implicitly locked on x86.)
  \item[Compare-and-swap] \code{CAS(\&x, expected, new)}: if \code{x ==
    expected}, write \code{new} --- all at once --- and report success.
    The universal building block: any read-modify-write can be written as
    a CAS loop (\code{lock cmpxchg}).
  \item[Fetch-and-add] \code{old = fetch\_add(\&x, n)} (\code{lock xadd}):
    counters, and the ticket lock below.
\end{description}

C11 exposes them in \code{<stdatomic.h>}; GCC also has builtins such as
\code{\_\_sync\_lock\_test\_and\_set}, which SWEB uses.

\section{Spinlocks}

\begin{lstlisting}
void spin_lock(atomic_int *l) {
  for (;;) {
    if (!atomic_exchange(l, 1)) return;   /* test-and-set: got it */
    while (atomic_load(l)) {}             /* test: wait with plain reads */
  }
}
void spin_unlock(atomic_int *l) { atomic_store(l, 0); }
\end{lstlisting}

Waiting with plain reads (\emph{test-and-test-and-set}) matters: every
\code{xchg} needs exclusive ownership of the cache line, so a crowd of
cores hammering \code{xchg} keeps moving the line between them and slows
down even the lock holder. Reads can share the line until the holder's
unlock invalidates it.

A plain spinlock is unfair --- whoever happens to win the race gets it.
The \textbf{ticket lock} is FIFO: \code{my = fetch\_add(\&next, 1); while
(serving != my) \{\}}; unlock increments \code{serving}. Like a
deli-counter ticket.

\section{Spin or sleep?}

Spinning wastes the CPU for the whole wait. That is fine when the lock is
held for a few instructions and the holder is running on another core.
It is disastrous when the holder is \emph{not running} --- on a single
CPU the spinner burns its entire time slice while the holder cannot make
progress. Hence the three options:

\begin{enumerate}
  \item spin (short critical sections, multi-core, holder never sleeps);
  \item spin, but \code{yield} after a few tries (gives the holder a chance
    to run);
  \item \textbf{sleep}: put yourself on a wait queue, block, and let the
    unlocker wake you. Costs two context switches, wastes no CPU. Linux's
    \code{futex} makes the uncontended case a single atomic instruction in
    user space and enters the kernel only to sleep and wake.
\end{enumerate}

\section{Memory ordering in one paragraph}

Atomic operations also order the \emph{other} memory accesses around
them. Lock acquisition needs \emph{acquire} semantics (nothing from the
critical section may move before it), unlock needs \emph{release}
semantics (nothing may move after it). A release store paired with an
acquire load that reads it creates a happens-before edge (chapter~3).
C11's default, \code{memory\_order\_seq\_cst}, is the strongest and
simplest; use it unless you know exactly why a weaker order suffices.

\begin{swebbox}
\file{common/source/kernel/}: \code{SpinLock} tries
\code{ArchThreads::testSetLock} (an \code{xchg}); if the lock is taken it
records itself as waiting and loops \emph{test-and-set + yield} ---
spinning without yielding would be pointless on SWEB's single CPU.
\code{Mutex} sleeps instead: on contention it enqueues the thread in its
waiters list and calls \code{Scheduler::sleep()}; \code{release} pops a
waiter and wakes it. Both inherit from \code{Lock}, which keeps the
per-thread lists used for deadlock checks (chapter~6). User-space locks
in SWEB (\code{pthread\_spin\_*}, and mutexes that sleep only through a
small system call) are part of P1's elective.
\end{swebbox}

\begin{keybox}
\begin{itemize}
  \item Requirements: mutual exclusion, progress, bounded waiting.
  \item Disabling interrupts works only on one CPU and only in the kernel.
  \item Peterson needs sequential consistency; store buffers break it
    without fences.
  \item Locks are built from atomic read-modify-write instructions
    (xchg, CAS, fetch-and-add).
  \item Spin briefly, otherwise sleep; on one CPU never spin without
    yielding.
  \item Lock = acquire, unlock = release: that is what makes the data
    written inside visible to the next holder.
\end{itemize}
\end{keybox}
